\documentclass[a4paper,11pt]{article}

% ---------------- Packages ----------------
\usepackage{amsmath,amssymb}
\usepackage{mathpartir}
\usepackage{geometry}
\geometry{margin=1in}

% ---------------- Macros ----------------
\newcommand{\Int}{\texttt{Int}}
\newcommand{\Bool}{\texttt{Bool}}
\newcommand{\ctx}{\Gamma}
\newcommand{\vctx}{\Sigma}

\newcommand{\typed}[3]{#1 ; #2 \vdash #3}
\newcommand{\checkstmt}[3]{#1 ; #2 \vdash #3\ \checkmark}

% ---------------- Document ----------------
\begin{document}

\title{Typing Rules}
\author{}
\date{}
\maketitle

\section*{Types}
\[
\tau ::= \Int \mid \Bool \mid V
\]

\section*{Typing Judgements}

\begin{itemize}
  \item Expression typing:
  \[
    \ctx ; \vctx \vdash e : \tau
  \]
  \item Statement typing:
  \[
    \ctx ; \vctx \vdash s \;\checkmark (\checkmark indicates \ well-typed)
  \]
\end{itemize}

\section*{Literals}

\begin{mathpar}
\inferrule*[right=T-Num]
{ }
{ \ctx ; \vctx \vdash n : \Int }

\inferrule*[right=T-Bool]
{ }
{ \ctx ; \vctx \vdash b : \Bool }
\end{mathpar}

\section*{Variables}

\begin{mathpar}
\inferrule*[right=T-Var]
{ x : \tau \in \ctx }
{ \ctx ; \vctx \vdash x : \tau }
\end{mathpar}

\section*{Binary Operations}

\subsection*{Arithmetic}

\begin{mathpar}
\inferrule*[right=T-Arith]
{
  \ctx ; \vctx \vdash e_1 : \Int \\
  \ctx ; \vctx \vdash e_2 : \Int
}
{
  \ctx ; \vctx \vdash e_1 \oplus e_2 : \Int
}
\quad
(\oplus \in \{+, -, *, /\})
\end{mathpar}

\subsection*{Comparison}

\begin{mathpar}
\inferrule*[right=T-Comp]
{
  \ctx ; \vctx \vdash e_1 : \Int \\
  \ctx ; \vctx \vdash e_2 : \Int
}
{
  \ctx ; \vctx \vdash e_1 \bowtie e_2 : \Bool
}
\quad
(\bowtie \in \{<, >\})
\end{mathpar}

\subsection*{Logical}

\begin{mathpar}
\inferrule*[right=T-Logic]
{
  \ctx ; \vctx \vdash e_1 : \Bool \\
  \ctx ; \vctx \vdash e_2 : \Bool
}
{
  \ctx ; \vctx \vdash e_1 \odot e_2 : \Bool
}
\quad
(\odot \in \{\land, \lor\})
\end{mathpar}

\subsection*{Equality}

\begin{mathpar}
\inferrule*[right=T-Eq]
{
  \ctx ; \vctx \vdash e_1 : \tau \\
  \ctx ; \vctx \vdash e_2 : \tau
}
{
  \ctx ; \vctx \vdash e_1 == e_2 : \Bool
}
\end{mathpar}

\section*{Unary Operations}

\begin{mathpar}
\inferrule*[right=T-Not]
{
  \ctx ; \vctx \vdash e : \Bool
}
{
  \ctx ; \vctx \vdash !e : \Bool
}
\end{mathpar}

\section*{Variants}

\subsection*{Variant Declaration}

\begin{mathpar}
\inferrule*[right=T-VariantDecl]
{
  V \notin \vctx
}
{
  \ctx ; \vctx \vdash
  \texttt{variant } V =
  \{ C_i(\tau_i) \}
  \Rightarrow
  \vctx \cup \{ C_i : \tau_i \to V \}
}
\end{mathpar}

\subsection*{Variant Constructor Use}

\begin{mathpar}
\inferrule*[right=T-VariantConstr]
{
  \ctx ; \vctx \vdash e_j : \tau_j
  \\
  \vctx(C) = (\tau_1,\dots,\tau_n) \to V
}
{
  \ctx ; \vctx \vdash C(e_1,\dots,e_n) : V
}
\end{mathpar}

\section*{Declarations and Assignment}

\subsection*{Variable Declaration}

\begin{mathpar}
\inferrule*[right=T-Decl]
{
  x \notin \ctx
}
{
  \ctx ; \vctx \vdash \texttt{decl } x : \tau
  \Rightarrow
  \ctx \cup \{ x : \tau \}
}
\end{mathpar}

\subsection*{Assignment}

\begin{mathpar}
\inferrule*[right=T-Assign]
{
  x : \tau \in \ctx \\
  \ctx ; \vctx \vdash e : \tau
}
{
  \checkstmt{\ctx}{\vctx}{x := e}
}
\end{mathpar}

\section*{Control Flow}

\subsection*{If}

\begin{mathpar}
\inferrule*[right=T-If]
{
  \ctx ; \vctx \vdash e : \Bool \\
  \forall s \in S.\; \checkstmt{\ctx}{\vctx}{s}
}
{
  \checkstmt{\ctx}{\vctx}{\texttt{if } e \texttt{ then } S}
}
\end{mathpar}

\subsection*{If--Else}

\begin{mathpar}
\inferrule*[right=T-IfElse]
{
  \ctx ; \vctx \vdash e : \Bool \\
  \forall s \in S_1.\; \checkstmt{\ctx}{\vctx}{s} \\
  \forall s \in S_2.\; \checkstmt{\ctx}{\vctx}{s}
}
{
  \checkstmt{\ctx}{\vctx}{
    \texttt{if } e \texttt{ then } S_1 \texttt{ else } S_2
  }
}
\end{mathpar}

\section*{Print}

\begin{mathpar}
\inferrule*[right=T-Print]
{
  \ctx ; \vctx \vdash e : \tau
}
{
  \checkstmt{\ctx}{\vctx}{\texttt{print } e}
}
\end{mathpar}

\section*{Program}

\begin{mathpar}
\inferrule*[right=T-Program]
{
  \forall s \in P.\; \checkstmt{\ctx}{\vctx}{s}
}
{
  \ctx ; \vctx \vdash P \;\checkmark
}
\end{mathpar}

\end{document}
